\bookchapter{Preface}{ch:00-preface}

\begin{chapterintro}

I use apps every day, but I have no idea what happens after I tap a button.

\end{chapterintro}

If that sounds like you, this short book is for you. It answers one question: \emph{what happens between tapping a button and seeing the result?}

\section{Who This Book Is For}\label{00-preface:who-this-book-is-for}

You use apps and websites every day, but you have never built one. You may know basic Python, which helps with one or two short examples, but you don't need to know anything about the web, servers, or networks. This book explains; it doesn't teach you to build. When you finish, you will have a clear picture of how any app works, ready for the books that teach the building.

\section{How This Book Is Organized}\label{00-preface:how-this-book-is-organized}

\begin{itemize}
\tightlist
\item
  Part I, \emph{Understand} (Chapters 1--4), builds the picture. Chapter 1 shows the two sides of every app. Chapter 2 looks at the side on your screen, Chapter 3 at the side you can't see, and Chapter 4 follows one tap all the way through both.
\item
  Part II, \emph{What Next} (Chapters 5--6), asks what can go wrong, what makes an app slow, how it breaks, and what it costs, then points to what to learn next.
\item
  The Glossary defines every new term in a line or two.
\end{itemize}

\section{How to Use This Book}\label{00-preface:how-to-use-this-book}

Read it in order; it is short, and each chapter builds on the one before. Chapters 1 to 4 each end with a \emph{Try It} section that you do in your own web browser on a computer. Nothing needs installing. Do them: seeing the requests fly past makes the ideas stick.

\section{Conventions Used in This Book}\label{00-preface:conventions-used-in-this-book}

\begin{itemize}
\tightlist
\item
  Bold type marks a new term the first time it is explained. Every such term is also in the Glossary.
\item
  \texttt{Constant\ width} marks code, commands, and anything you would type exactly.
\item
  Code blocks show short examples and their real output. Commands are shown without a prompt.
\item
  Three kinds of callout appear in the text:
\end{itemize}

\begin{callout}{tip}

A suggestion that saves time or trouble.

\end{callout}

\begin{callout}{note}

A general remark, or something to be aware of.

\end{callout}

\begin{callout}{warning}

A mistake that can cost you data, money, or security.

\end{callout}
